% GENERATED FILE. Edit canonical lessons, then run npm run generate:books.
% canonical-source-hash: fnv1a64:eea35edc6ab07100
% canonical-lessons: RU-C64-yellow, RU-C64-beautiful, RU-C64-good, RU-C64-bad, RU-C64-new

\chapter{Yellow, Beautiful, Good, Bad, New}
\label{ch:ru-yellow-beautiful-good-bad-new}
\hlchaptermodality{64}

\begin{chapteropening}
\textbf{By the end of this chapter:} \emph{I can say yellow, beautiful, good, bad and new.}

Complete the last lesson of chapter 64: I can say yellow, beautiful, good, bad and new.
\end{chapteropening}

\section[\texorpdfstring{\ru{жёлтый} (zhyóltyy)}{жёлтый (zhyóltyy)}]{\ru{жёлтый} (zhyóltyy) — yellow}

\label{lesson:RU-C64-yellow}



\begin{quote}
Before the new one: say the Russian for blue, then the Russian for green.
\end{quote}

\subsection*{You'll want to know: \ru{жёлтый}}
\textbf{\ru{жёлтый}} (\emph{zhyóltyy}) — \textquotedblleft{}yellow\textquotedblright{}.

\begin{cousinweb}
\textbf{\ru{жёлтый}} $\leftarrow$ Proto-Slavic \emph{*žĭltŭ} — inherited. Vasmer relates it to Lithuanian \emph{geltas}, \textquotedblleft{}yellow\textquotedblright{}, Greek \emph{chol\'{\={e}}}, \textquotedblleft{}bile\textquotedblright{}, and Latin \emph{fel}, \textquotedblleft{}bile\textquotedblright{}: bile was named for its colour. \emph{\ru{желчь}}, \textquotedblleft{}bile\textquotedblright{}, is its Russian twin.
\end{cousinweb}

\begin{grammarlens}[title={What you've built}]
One more word for this chapter.
\end{grammarlens}

\subsection*{Your turn}
\begin{itemize}
\raggedright
  \item \emph{Say it:} \emph{zhyóltyy}
  \item \emph{Say it:} \emph{zhyóltyy}, once more
  \item \emph{Say it:} say \emph{zhyóltyy}
  \item \emph{Recall:} say the Russian for blue, then the Russian for green, then say \emph{zhyóltyy} again
\end{itemize}

\begin{tcolorbox}[breakable,colback=teal!4,colframe=teal!35!black,title={Before you move on}]
What does \ru{жёлтый} mean? (\textquotedblleft{}Yellow\textquotedblright{}.) Say it once more.
\end{tcolorbox}
\section[\texorpdfstring{\ru{красивый} (krasívyy)}{красивый (krasívyy)}]{\ru{красивый} (krasívyy) — beautiful}

\label{lesson:RU-C64-beautiful}



\begin{quote}
Before the new one: say the Russian for green, then the Russian for yellow.
\end{quote}

\subsection*{You'll want to know: \ru{красивый}}
\textbf{\ru{красивый}} (\emph{krasívyy}) — \textquotedblleft{}beautiful\textquotedblright{}.

\begin{cousinweb}
\textbf{\ru{красивый}} is \emph{\ru{краса}}, \textquotedblleft{}beauty\textquotedblright{}, with \emph{-\ru{ивый}}, \textquotedblleft{}full of\textquotedblright{} (Vasmer); \emph{\ru{красный}}, \textquotedblleft{}red\textquotedblright{}, once \textquotedblleft{}beautiful\textquotedblright{}, comes from the same word.
\end{cousinweb}

\begin{grammarlens}[title={What you've built}]
One more word for this chapter.
\end{grammarlens}

\subsection*{Your turn}
\begin{itemize}
\raggedright
  \item \emph{Say it:} \emph{krasívyy}
  \item \emph{Say it:} \emph{krasívyy}, once more
  \item \emph{Say it:} say \emph{krasívyy}
  \item \emph{Recall:} say the Russian for green, then the Russian for yellow, then say \emph{krasívyy} again
\end{itemize}

\begin{tcolorbox}[breakable,colback=teal!4,colframe=teal!35!black,title={Before you move on}]
What does \ru{красивый} mean? (\textquotedblleft{}Beautiful\textquotedblright{}.) Say it once more.
\end{tcolorbox}
\section[\texorpdfstring{\ru{хороший} (khoróshiy)}{хороший (khoróshiy)}]{\ru{хороший} (khoróshiy) — good}

\label{lesson:RU-C64-good}



\begin{quote}
Before the new one: say the Russian for yellow, then the Russian for beautiful.
\end{quote}

\subsection*{You'll want to know: \ru{хороший}}
\textbf{\ru{хороший}} (\emph{khoróshiy}) — \textquotedblleft{}good\textquotedblright{}.

\begin{cousinweb}
\textbf{\ru{хороший}} — recorded from the 1200s. Vasmer thinks it most likely a short form of \emph{\ru{хоробрый}}, \textquotedblleft{}brave\textquotedblright{}; a link with \emph{\ru{хоронить}}, \textquotedblleft{}to guard\textquotedblright{}, he finds less likely, and one with the sun god Khors, which Abaev later supported, he rejects.
\end{cousinweb}

\begin{grammarlens}[title={What you've built}]
One more word for this chapter.
\end{grammarlens}

\subsection*{Your turn}
\begin{itemize}
\raggedright
  \item \emph{Say it:} \emph{khoróshiy}
  \item \emph{Say it:} \emph{khoróshiy}, once more
  \item \emph{Say it:} say \emph{khoróshiy}
  \item \emph{Recall:} say the Russian for yellow, then the Russian for beautiful, then say \emph{khoróshiy} again
\end{itemize}

\begin{tcolorbox}[breakable,colback=teal!4,colframe=teal!35!black,title={Before you move on}]
What does \ru{хороший} mean? (\textquotedblleft{}Good\textquotedblright{}.) Say it once more.
\end{tcolorbox}
\section[\texorpdfstring{\ru{плохой} (plokhóy)}{плохой (plokhóy)}]{\ru{плохой} (plokhóy) — bad}

\label{lesson:RU-C64-bad}



\begin{quote}
Before the new one: say the Russian for beautiful, then the Russian for good.
\end{quote}

\subsection*{You'll want to know: \ru{плохой}}
\textbf{\ru{плохой}} (\emph{plokhóy}) — \textquotedblleft{}bad\textquotedblright{}. The opposite of \textbf{\ru{хороший}}.

\begin{cousinweb}
\textbf{\ru{плохой}} $\leftarrow$ Old Russian \emph{ploxŭ}, \textquotedblleft{}bad\textquotedblright{}; Ukrainian \emph{\ru{плохий}} means \textquotedblleft{}meek, quiet\textquotedblright{}, and Polish \emph{płochy}, \textquotedblleft{}timid\textquotedblright{}. Vasmer links it with \emph{\ru{переполох}}, \textquotedblleft{}a fright\textquotedblright{}, while others relate it to \emph{\ru{плоский}}, \textquotedblleft{}flat\textquotedblright{}, so the origin is not settled.
\end{cousinweb}

\begin{grammarlens}[title={What you've built}]
One more word for this chapter.
\end{grammarlens}

\subsection*{Your turn}
\begin{itemize}
\raggedright
  \item \emph{Say it:} \emph{plokhóy}
  \item \emph{Say it:} \emph{plokhóy}, once more
  \item \emph{Say it:} say \emph{plokhóy}
  \item \emph{Recall:} say the Russian for beautiful, then the Russian for good, then say \emph{plokhóy} again
\end{itemize}

\begin{tcolorbox}[breakable,colback=teal!4,colframe=teal!35!black,title={Before you move on}]
What does \ru{плохой} mean? (\textquotedblleft{}Bad\textquotedblright{}.) Say it once more.
\end{tcolorbox}
\section[\texorpdfstring{\ru{новый} (nóvyy)}{новый (nóvyy)}]{\ru{новый} (nóvyy) — new}

\label{lesson:RU-C64-new}



\begin{quote}
Before the new one: say the Russian for good, then the Russian for bad.
\end{quote}

\subsection*{You'll want to know: \ru{новый}}
\textbf{\ru{новый}} (\emph{nóvyy}) — \textquotedblleft{}new\textquotedblright{}.

\begin{cousinweb}
\textbf{\ru{новый}} $\leftarrow$ Proto-Slavic \emph{*novŭ} — inherited from PIE \emph{*néwos}, \textquotedblleft{}new\textquotedblright{}. Vasmer lists Sanskrit \emph{náva-}, Greek \emph{néos}, Latin \emph{novus} and Gothic \emph{niujis}: English \emph{new} is a cousin.
\end{cousinweb}

\begin{grammarlens}[title={What you've built}]
One more word for this chapter.
\end{grammarlens}

\subsection*{Your turn}
\begin{itemize}
\raggedright
  \item \emph{Say it:} \emph{nóvyy}
  \item \emph{Say it:} \emph{nóvyy}, once more
  \item \emph{Say it:} say \emph{nóvyy}
  \item \emph{Recall:} say the Russian for good, then the Russian for bad, then say \emph{nóvyy} again
\end{itemize}

\begin{tcolorbox}[breakable,colback=teal!4,colframe=teal!35!black,title={Before you move on}]
What does \ru{новый} mean? (\textquotedblleft{}New\textquotedblright{}.) Say it once more.
\end{tcolorbox}
